\documentclass[10pt,a4paper]{article}
\usepackage[margin=1.75cm,top=1.45cm,bottom=1.45cm]{geometry}
\usepackage[T1]{fontenc}
\usepackage[utf8]{inputenc}
\usepackage{lmodern}
\usepackage{microtype}
\usepackage{graphicx}
\usepackage{float}
\usepackage{xurl}
\usepackage{booktabs}
\usepackage{longtable}
\usepackage{array}
\usepackage{xcolor}
\usepackage[hidelinks]{hyperref}
\usepackage[numbers,sort&compress]{natbib}
\usepackage{caption}
\captionsetup{font=small,labelfont=bf,skip=4pt}
\usepackage{titlesec}
\titlespacing*{\section}{0pt}{8pt}{3pt}
\titlespacing*{\subsection}{0pt}{6pt}{2pt}
\titleformat*{\section}{\normalsize\bfseries}
\titleformat*{\subsection}{\normalsize\bfseries}
\setlength{\parskip}{2pt}
\setlength{\parindent}{10pt}
\setlength{\textfloatsep}{8pt}
\setlength{\intextsep}{6pt}
\setlength{\abovecaptionskip}{3pt}
\graphicspath{{../out/figures/}}
\input{../out/macros.tex}

\title{\vspace{-14pt}\Large Watching AI model documentation:\\ what changes, what goes quiet, and what the watcher misses\vspace{-4pt}}
\author{\normalsize AI-generated report (cardtrack-report skill)\thanks{Written by an AI agent following the
cardtrack-report skill; see the disclaimer on systemcards.org/analysis for the model and skill version. Every number,
table and figure is generated by \texttt{analyze.py} from a read-only snapshot of the cardtrack database. Version
pairs were classified by AI sub-agents and spot-checked by the generating agent. Appendix~\ref{app:changes} lists
every substantive change with its URL and version ids.}}
\date{\normalsize Data snapshot \snapshotDatePretty}

\begin{document}
\maketitle
\vspace{-18pt}

\begin{abstract}
\noindent
cardtrack is a daily tracker of first-party model and system cards, access policies and independent
evaluations of frontier AI models, with additions proposed by a large language model (LLM) agent and
audited by a human (\href{https://systemcards.org}{systemcards.org}). We analyse its first \trackingWeeks{}
weeks: \nTracked{} documents from \nPublishers{} publishers, \nVersions{} stored versions and \nChecks{} link
and content checks. Post-publication edits are real and worth tracking: \nPubSubstantive{} publisher-side
version changes were substantive, among them \nFlagCorrection{} corrections of a stated number or fact (one propagated through
\notablePropagatedMax{} OpenAI documents), scope changes to access policies, and score edits and removals
with no change log; \nDocsSubstantive{} documents (\shareDocsSubstantive\%) changed substantively within the
window. The cost of finding them is that \sharePubNonchange\% of detected changes were page furniture or
extraction artefacts, and apparent deletions must be checked against the raw capture before they are
believed. Hard loss is rare: \docsEverNotFound{} of \nDocsWithValidChecks{} documents ever returned a 404,
\nDeadWord{} of them for good, too few to estimate a link-rot rate over this window. The quieter failure is
a page that keeps answering HTTP 200: \nSoftGoneWord{} pages served a redirect stub for at least
\softGoneDaysMin{} days, visible only to content fingerprinting until the publisher added real redirects. A
host-side fault cost \shareOutageRuns\% of the daily runs their checks, none since \lastOutagePretty{}
(Appendix~\ref{app:runs}), so the change counts above are lower bounds
and their detection dates run late.
\end{abstract}

\section{Setting and data}
Each daily run probes every tracked URL (link check), re-fetches a 15\,\% rotation of documents and compares
a text fingerprint that ignores known page furniture such as download counters, sidebars and footers
(content check), diffs publisher index pages for new links, and then lets a sandboxed LLM agent propose
additions through a deterministic validator. A supervised backfill on \firstSeenMinPretty{} seeded
\nBackfill{} documents; the snapshot covers \nRuns{} runs from \firstRunPretty{} to \snapshotDatePretty,
during which \nPostBackfill{} more were discovered. The only continuous public log
of edits to such documents we know of is the Midas Project's hand-curated Watchtower \cite{midas2026}; this
study is fingerprint-based and model-centred.

\section{What changes}
\begin{figure}[H]
\centering
\includegraphics[width=\linewidth]{fig_changes}
\caption{(a) Classification of stored version pairs by what the text diff contained (groups with ten or
more pairs; the ``before filter fix'' row lacks the versions later deleted; see text). (b) Kaplan--Meier estimate, that is the share of documents with at least one substantive
publisher-side change by day $t$ adjusted for documents tracked fewer than $t$ days, by the document's
current format, drawn while at least \kmMinAtRisk{} documents remain at risk; changes are dated at detection, and backfilled documents were already months old on day 0.}
\label{fig:changes}
\end{figure}
Of the \nPairs{} consecutive version pairs in the database, \nMigrationPairs{} (\shareMigrationPairs\%)
are the tracker's own doing, so the diff compares two captures, not two revisions: canonical URLs
migrated on \filterRecomputePretty{} from announcement pages to full PDFs, and first captures of a redirect
target after a recorded move. The remaining \nPubPairs{} publisher-side
pairs were classified by reading the text diffs (Figure~\ref{fig:changes}a): \sharePubNonchange\% were
Hugging Face download counters, evaluation-widget rows, footers or PDF re-extraction artefacts,
\sharePubSubstantive\% changed the document's content and \sharePubMajor\% were major edits. On
\filterRecomputePretty{} the fingerprint was also recomputed to ignore known furniture. This cut the
per-check change-flag rate from \shareFpChangedPreFilter\% to \shareFpChangedPostFilter\% while the number
of substantive pairs found did not fall (\subPre{} before, \subPost{} after, over a longer window), so
precision per flag rose from
about \precisionPre\% to \precisionPost\%. The \emph{before} side is incomplete: \nVersionsMissing{} versions that
the changelog records as written are no longer in the database, \nVersionsMissingPreFilter{} of them from
before the recompute; \nClassifiedPurgedWord{} of them are the footnote purge described below and the cause of the
rest is not recorded, so pairs of every kind are missing there. The remaining leak is evaluation-widget rows that end in a score,
which the bullet-anchored patterns do not match. In document terms, \nDocsSubstantive{} of \nTracked{}
tracked documents (\shareDocsSubstantive\%) changed substantively during the window, an estimated
\kmHtmlThirty\% of HTML and \kmPdfThirty\% of PDF documents within 30 days of tracking; the two curves
in Figure~\ref{fig:changes}b are not distinguishable at this sample size.

The substantive changes are the reason to run a tracker at all (Table~\ref{tab:notable};
Appendix~\ref{app:changes} lists all \nAppendixRows{} with URLs). Of the publisher-side pairs,
\nFlagCorrection{} corrected a stated number or fact, \nFlagScore{} touched an evaluation number and \nFlagSafety{}
touched a safety, risk or mitigation section. Three patterns recur (Appendix~\ref{app:changes} ids in
brackets). First, some corrections come with dated change logs, mostly at OpenAI: one corrected number
propagated with its change log into \notablePropagatedMax{} OpenAI documents that quote it as a baseline
(\exPropagated), and the GPT-6 Astra system card gained an appendix covering GPT-6 Sol and Luna
and corrected Astra's HealthBench results after an evaluation misconfiguration, with dated entries
(\exAstraAppendix), then, with another dated entry, a second appendix on ``dots'', always-on agents built on
Astra, with prompt-injection red-teaming and monitor-evasion results (\exAstraDots); xAI (\exGrokLog) and Anthropic (\exAnthropicLog) each added one dated log. Second, edits
without a change log exist: an NVIDIA card revised a tool-calling score downward and dropped its research-only
use restriction with no note (\exVoiceChat).
A date label does not settle when an edit was made either: METR replaced a one-line pointer on its red-teaming of Anthropic's internal agent monitoring with the findings
themselves, including routes to disabling monitors, under a heading dated months before the tracker's first
capture, although every earlier successful content check found the page unchanged (\exMetrMonitoring). Third,
documents change scope: OpenAI rebuilt its Daybreak cyber-access tiers to add GPT-6 Sol, Luna and Astra and to
require extra approval for GPT-5.6-Cyber (\exDaybreakTiers); Anthropic's cyber-safeguards page now excludes
Opus 5.5 and Sonnet 5.5 (\exCyberScope); NVIDIA's Cosmos 3 cards were each cut back to a single model
(\exCosmos); and inclusionAI added a ``training content summary'' linking a public training-data statement
to several of its cards (\exTrainingSummary). A caution the other way: the previous edition found that apparent deletions of
Anthropic footnotes were extractor artefacts. The publisher had wrapped its footnotes in an element whose class
name contains ``footer'', which the boilerplate remover discards, while the raw HTML still held every
footnote. The extractor has since been fixed and \nClassifiedPurgedWord{} classified versions, the artefacts among them, were deleted
from the database.\extractorArtefactSentence{} A fingerprint tracker cannot tell disclosed from
undisclosed edits, or edits from extraction bugs, but it produces the diff that lets a human tell.

\begin{table}[!htb]
\centering\scriptsize
\input{../out/tables/notable_report.tex}
\caption{Publisher-side changes selected by rule from Appendix~\ref{app:changes} (edits the classifier described as undisclosed, dated
corrections, pages whose content moved behind an HTTP 200, licence changes, then the largest major edit),
one row per distinct change, capped per kind. Summaries are the classifier's; URLs and version ids are in
Appendix~\ref{app:changes}.}
\label{tab:notable}
\end{table}

\section{What goes quiet}
Excluding host-side outage runs and the \nTrackerErrorsValidRuns{} probes that never reached a server leaves
\nValidLinkChecks{} valid link checks over \nDocsWithValidChecks{} documents (\docYearsObserved{}
document-years, including documents later removed from the corpus), of which \shareLinkOk\% succeeded.
Hard loss is rare: \docsEverNotFound{} documents ever returned a 404 (\notfoundDocsPerDocYear{} per
document-year). The monitor declared \nDeadWord{} of them dead after three consecutive 404s, a FAR.AI
evaluation of the cyber safeguards of Alibaba's Qoder agent (Qwen3.8-Max); the other was a slug typo that the publisher fixed with a
redirect. \docsEverNotFoundWordCap{} cases cannot give a reliable rate over \trackingWeeks{} weeks, but they are not out of line with
published baselines for young pages: Pew found 8\,\% of pages under one year old inaccessible
\cite{pew2024}, deep links in the Wayback corpus have a median lifetime of 1.3 years \cite{garg2024}, and 1
to 4\,\% of web references in scholarly articles from 2012 were already rotten when checked
\cite{klein2014}. Moves are more common: \nMovedEvents{} documents
moved via permanent redirect (\movesPerDocYear{} per document-year: site restructures, a trailing slash, the
typo fix and the soft-gone pages below). Bot blocking is the larger nuisance: \docsEverBlocked{} documents
were refused with HTTP 403 at least once, \topBlockedHostDocs{} of them OpenAI
documents on \topBlockedHost. For most of the window the blocks were intermittent (overall a median of \blockedPersistenceMedian\% of
later checks), consistent with rate-based bot detection, but in the last week most of the blocked OpenAI pages
were refused in most runs; whether that reflects the site or the tracker's move to a server host at the end of September is not known, and block
rates depend on the fetch client's fingerprint as much as on the site \cite{gundelach2026}. This matters
because OpenAI is also the publisher whose documents carried most of the dated corrections above. The
quietest failures were invisible to the link checker: \nSoftGoneWord{} pages kept returning HTTP 200 while
their body became a two-line client-side redirect stub (\texttt{meta refresh} plus a canonical link to the
new location), and only the content fingerprint noticed. For at least \softGoneDaysMin{} days the link
checker saw nothing, until the publisher replaced the stubs with real permanent redirects on
\softGoneResolvedPretty{} and the monitor marked them moved (Table~\ref{tab:softgone}). This is the reference
rot of Zittrain et al.\ \cite{zittrain2014}, the URL resolving while the content is gone, although here it
was temporary because the publisher had moved rather than removed the content.

\section{Discovery is fast}
For the \nLagLive{} documents published after the backfill, the median delay from publication to first
tracking was \lagMedianDays{} days and \lagWithinSevenDays\% were tracked within a week. Leads from index
diffs and from the monitor's own candidate queue were tracked within two days (median
\lagLeadFastRange{} days); free-form agent search took \lagLeadAgentSearch{} days and manual submissions
\lagLeadManual{} days.

\section{Limitations}
The window is \trackingWeeks{} weeks on two hosts in turn (a laptop, then a server) with one fetch client: enough to measure short-run edit
rates, not link rot, whose baselines are annual; block rates in the literature range from under 1\,\% for
browser-like clients to 15\,\% for headless ones \cite{gundelach2026}. Coverage was uneven: \shareOutageRuns\%
of the daily runs lost their checks to a host-side fault (Appendix~\ref{app:runs}), which stretched the content
re-check interval's tail to \fpGapPNineZeroDays{} days at the 90th percentile against an intended
\fpIntendedGapDays{} days, and content changes are observed only when a document's turn in the 15\,\% rotation comes
up, so change dates lag true edit dates by up to a re-check interval and edits reverted between checks are
missed. The instrument changed mid-window (URL migrations and the fingerprint recompute on
\filterRecomputePretty). The classification was done by LLM sub-agents with a written taxonomy (\nConfHigh{}
pairs at high confidence, \nConfMedium{} at medium, \nConfLow{} at low); it was spot-checked, not double-coded, categories at the
boundary of ``minor'' and ``metadata'' are soft, and text extraction can drop page elements, so apparent
deletions need a check against the raw capture. The corpus is curated by an allow-list and an agent, and the
literature baselines cover different populations.

\section{Future work}
On the science: a longer horizon allows a hazard model of edit rates by publisher and document type;
aligning detected edits with publishers' change logs and with Watchtower's disclosed and undisclosed tags
would estimate the disclosed fraction; classifying every minted version in the pipeline would turn the
appendix into a live feed. On the instrument: re-fetch every document daily instead of a 15\,\% rotation
(the whole corpus is under 400\,MB and the link probe already touches every URL), so edits are dated to the
day; treat a body that collapses to a redirect stub as a soft 404 and re-point moved documents; extend the footnote
pre-clean to the newer footnote markup, stop skipping tab panels, and diff the raw capture alongside the extracted text so extractor regressions
are never read as deletions; extend the furniture patterns to widget rows that end in a
score.

\bibliographystyle{unsrtnat}
\scriptsize
\setlength{\bibsep}{0pt}
\bibliography{refs}

\clearpage
\appendix
\normalsize
\section{Every substantive publisher-side change}
\label{app:changes}
Table~\ref{tab:appendix} lists all \nAppendixRows{} publisher-side version pairs classified as a content
change (minor, major, or replaced), in order of detection, with the canonical URL and the stored version ids
so the diff can be reproduced from the database (\texttt{document\_versions} rows, text under
\texttt{data/text/}). Detection dates lag the edit by up to one re-check interval. Table~\ref{tab:softgone}
lists the pages that went quiet behind an HTTP 200 and Table~\ref{tab:moved} the recorded moves.
\begin{center}\scriptsize
\input{../out/tables/appendix_changes.tex}
\end{center}
\begin{center}\scriptsize
\input{../out/tables/appendix_soft_gone.tex}
\end{center}
\begin{center}\scriptsize
\input{../out/tables/appendix_moved.tex}
\end{center}

\section{Tracker reliability}
\label{app:runs}
Of the \nRuns{} runs, \nOutageRuns{} recorded a fetch error for every URL, the longest streak lasting
\nOutageStreakDays{} consecutive runs (Figure~\ref{fig:heatmap}); the last was on \lastOutagePretty, and none of the
\nRunsSinceLastOutage{} runs since then was a full outage. \outageCauseSentence{} An outage run still continues to its agent phase, which can succeed
(\nOutageRunsWithAdds{} outage runs still added documents), but that day's link and content checks are lost. A further \nPartialRunsWord{} runs lost
part of their checks, the worst \worstPartialErrors{} of \worstPartialN{} probes on \worstPartialDate.
\agentPhaseSentence{} Run logs for the
other \nRunsAgentUnknown{} runs are not in the snapshot, so their agent phase is unknown (``no log'' in
Table~\ref{tab:runs}). Outages stretched the tail of the content re-check interval.
Table~\ref{tab:runs} gives every run.
\begin{figure}[H]
\centering
\includegraphics[width=0.9\linewidth]{fig_heatmap}
\caption{Link-check outcome of every tracked document (rows, grouped by publisher, ordered by date first
tracked; the largest publishers labelled where space allows) in every run (columns). Grey columns: every
check failed on the tracker's side; light blue: reachable; orange: HTTP 403; red: 404; violet: permanent
redirect; amber: fetch error in a healthy run; white: not yet tracked.}
\label{fig:heatmap}
\end{figure}
\begin{center}\scriptsize\renewcommand{\arraystretch}{0.88}
\input{../out/tables/runs_by_day.tex}
\end{center}
\end{document}
